\section{Centered cell budgets imply finite-cutoff control}\label{ex-centered-estimates}

This section proves the quantitative step used before taking any continuum limit. The estimates are conditional on a centered cell bound along one fixed finite-volume trajectory. They do not produce that bound from the renormalisation-group flow; the latter is a separate and substantially stronger input recorded in the blueprint. The point of the argument here is to show in detail how a bound on one observable centered in one cell controls arbitrary separated tuples, and how the same bound controls changes of lattice positions.

Write $a_N$ for the lattice spacing at mesh index $N$, and $\ell_{N,p}$ for the number of sites in one coordinate direction of the period-$p$ torus. Thus its physical linear period is $P_p=a_N\ell_{N,p}$, independent of $N$ when the schedule is embedded. A configuration of $q$ marked points is a point of $E_q=(\mathbb R^4)^q$. Let $\mathrm{Conf}_q$ be the set of configurations with no coincident points, and let $\mathrm{Gen}_q$ consist of configurations for which each coordinate projection is injective on the marked points. For period $p$, let $\mathrm{Cyc}_p$ be the set where every pairwise coordinate difference avoids $P_p\mathbb Z$, and let $\mathrm{NoHit}_p$ be the set where each pair is nonperiodic in at least one coordinate. Thus $\mathrm{Gen}_q\subset\mathrm{Conf}_q$, while $\mathrm{NoHit}_p$ permits coordinatewise ties. We use the sup norm on each spatial coordinate when discussing separations. A tuple is \emph{cyclically separated at scale $r$} if it lies in $\mathrm{Cyc}_p$ and, for each pair and coordinate, its coordinate difference has distance at least $r$ from every multiple of $P_p$. It is \emph{separated in some coordinate at scale $r$} if for each pair at least one coordinate has this property. No condition is imposed on tuples with coincident marked points when we assert a compact-set bound only away from collisions.

For a local observable $O$ (that is, one of the tokens in the fixed observable register of the finite-volume family) and a lattice site $x$, let $O_N(x)$ be its normalized finite-volume readout. For a finite list $\mathbf O=(O_0,\ldots,O_{q-1})$ and sites $\mathbf x=(x_0,\ldots,x_{q-1})$, write
\[
  F_{N,p,\mathbf O}(\mathbf x)
  =\mathbb E_{N,p}\!\left[\prod_{i=0}^{q-1}O_i{}_N(x_i)\right].
\]
The same notation is used for the lifted lattice function on $E_q$. The list is allowed to contain any fixed local observables, including reflected or conjugated observables. A centered cell of physical width comparable to $a_m$ is the centered link tiling at depth $m\leq N$, with its first chamber centered at the chosen site. Call a complex observable chamber-admissible if it is gauge invariant, measurable, uniformly bounded on the finite-volume configuration space, and depends only on link variables whose two endpoints lie in that chamber. Admissibility does not mean that the observable is pointwise nonnegative; reflection positivity is a property of the expectation and the reflected product.

\begin{definition}[Centered budget]\label{ex-centered-budget}
For an observable $O$, a nonnegative number $A$, and indices $m\leq N$, a centered budget $\mathsf{CB}(O;A,m,N,p)$ means: (i) the depth of $O$ is represented at volume $N$; (ii) at every torus site $x$, the readout $O_N(x)$ is chamber-admissible for the home chamber of the centered tiling of depth $m$ based at $x$; and (iii) its centered root test is at most $A$. The root test is the $M$th root of the absolute normalized expectation of the product of the observable copied, with the prescribed reflected/conjugated orientation, to all $M$ cells of that tiling. The reflected copies and the root test use the finite-volume reflection-positive transfer structure.
\end{definition}

For a centered tiling $\mathcal T$ with $M$ cells, write
\[
 \mathcal R_{\mathcal T}(F)
 :=\left|\mathbb E_{N,p}\left[\prod_{C\in\mathrm{Cell}(\mathcal T)} F^{(C)}\right]\right|^{1/M},
\]
where $F^{(C)}$ is the prescribed translated copy of $F$, reflected and conjugated on alternating chambers. This notation is used only when the observable is admissible in the home chamber and the normalized expectation is defined.

The quantitative assumption has two distinct orders of quantifiers. For every fixed observable $O$ and period index $p$, there is a tail of depths $m$ such that for each such $m$ one can choose $A=A(O,p,m)\geq0$ and a mesh tail on which $\mathsf{CB}(O;A,m,N,p)$ holds. In addition, for each $O$ there is a tail of depths $m$ such that one can choose $A=A(O,m)\geq0$ for which there is an outer period-index tail and, for each period in it, an inner mesh-index tail on which the budget holds. The constants may depend on the displayed fixed data. This is the precise centered-budget hypothesis used in the next result.

\begin{lemma}[Finite families share one budget]\label{ex-centered-finite}
Suppose each observable in a finite list has the fixed-period centered-budget property (respectively, the expanding-period property). Then, after increasing the depth tail, there is one nonnegative budget for the whole list and one common eventuality on which all the individual centered budgets hold. The same conclusion applies to any finite family of reflected or conjugated observables that belong to the register and have the corresponding centered-budget property.
\end{lemma}
\begin{proof}
For the empty list take budget zero. Otherwise, for each member $O_i$, choose its depth threshold $m_i$, and let $m_* = \max_i m_i$. Fix any $m\geq m_*$. At this particular depth, the hypotheses give a possibly depth-dependent budget $A_i(m)$ and an eventuality for $O_i$. Take $A(m)=\max_i A_i(m)$ and intersect the finitely many mesh tails (or the finitely many ordered $(N,p)$ tails). The result is a common eventuality and a single budget for the entire list at this same depth $m$. Since this works for every $m\geq m_*$, it proves the asserted eventual-in-depth statement. No monotonicity of the centered budgets as depth changes is used or claimed. There are finitely many orientations of each fixed observable: take all compositions of the four coordinate reflections, with or without conjugation. Whenever these are registered observables with the required property, include the entire finite family in the same list. This argument does not claim a bound uniform over all observables or over an uncountable test class.
\end{proof}

\begin{theorem}[Centered budgets give local moduli and compact bounds]\label{ex-centered-producers}
Assume the finite-volume Wilson family has the following properties, as well as the centered-budget property above. Its fixed schedule satisfies $a_N\to0$, $P_p\to\infty$ along the period schedule, and the nested lattice scales and centered tilings have the stated integer alignment; the finite-volume side is divisible into the $4s$ equal chambers used by the centered transfer word. The observable register is closed under the coordinate reflections and conjugation used below. Each registered observable has a fixed finite stencil, its readout is continuous and covariant under translations and coordinate reflections, and conjugation is implemented by the reflected token. Every coordinate cut has a positive seam and the associated continuous transfer representation: a compact metric slice space $X$, a finite Borel measure $\mu$, a continuous positive-semidefinite heat kernel $T$, and continuous kernels $B_i$ for local slab observables. If a cyclic word consists of bricks $B_i$ of common width $w$ and nonnegative integer free gaps $H_i$, with $qw+\sum_iH_i=L$, the unnormalized finite-volume expectation is the cyclic kernel integral
\[
 \operatorname{Cl}_\mu(K_1,\ldots,K_r)
 :=\int_{X^r}\prod_{j=1}^{r}K_j(x_j,x_{j+1})\,d\mu^r(x_1,\ldots,x_r),\qquad x_{r+1}=x_1,
\]
for the collapsed cyclic list of $q$ brick kernels and the individual unit heat steps, which has $r=q+\sum_i H_i$ kernel factors. The physical tour length is $L=qw+\sum_iH_i$. Division by the $L$-step partition trace $Z=\operatorname{Cl}_\mu(T,\ldots,T)>0$ gives the normalized Wilson expectation. The finite-volume representation identifies a change of one marked position by $l$ steps with replacing two free gaps $H_i,H_j$ by $H_i-l,H_j+l$. Finally, for every cut tiling $\mathcal T$, finite index set $I$, injection $c:I\to\mathrm{Cell}(\mathcal T)$, and home-chamber-admissible functions $F_i$, the planted inequality
\[
 \left|\mathbb E_{N,p}\!\left[\prod_{i\in I} F_i^{(c(i))}\right]\right|
 \leq \prod_{i\in I} \mathcal R_{\mathcal T}(F_i)
\]
holds. These are the precise locality, transfer, symmetry, reality, seam, and cell-packing properties used below. Then the following four conclusions hold.
\begin{enumerate}
\item For each fixed period $p$, each compact set $K\subset\mathrm{Cyc}_p$, each finite list $\mathbf O$, and every $\varepsilon>0$, there exist $\delta>0$ and a mesh threshold $N_0$ such that, for $N\geq N_0$, any two lattice tuples $\mathbf x,\mathbf y$ whose physical embeddings lie in $K$ and satisfy $|\mathbf x-\mathbf y|<\delta$ obey
\[
 |F_{N,p,\mathbf O}(\mathbf x)-F_{N,p,\mathbf O}(\mathbf y)|<\varepsilon.
\]
\item For each compact $K\subset\mathrm{Gen}_q$, each finite list $\mathbf O$, and every $\varepsilon>0$, there are $\delta>0$ and an outer period-index threshold such that for every period index in that tail the same mesh modulus holds eventually in the inner mesh index. The mesh threshold may depend on the period, while $\delta$ does not.
\item For each compact $K\subset\mathrm{Conf}_q$ and each finite list $\mathbf O$, there are $D<\infty$ and an outer period-index threshold such that for every period index in that tail the same finite-volume functions have absolute value at most $D$ eventually in the inner mesh index on lattice tuples embedded in $K$.
\item For each fixed period $p$, each compact $K\subset\mathrm{NoHit}_p$, and each finite list $\mathbf O$, there is $D_p<\infty$ such that $|F_{N,p,\mathbf O}(\mathbf x)|\leq D_p$ for all lattice tuples embedded in $K$, eventually in the mesh limit. This hypothesis permits coordinatewise ties and applies even when the points are not separated in every coordinate.
\end{enumerate}
The first modulus may depend on the fixed period; the second modulus is uniform in the outer period tail. The third bound has one constant in that tail, with a mesh threshold that may depend on the period. The fourth bound is only for a fixed period. No assertion extends a bound through point collisions.
\end{theorem}

\begin{proof}
If the compact set $K$ in any assertion is empty, that assertion is immediate. We therefore take $K$ nonempty when using compact minima below. For the modulus assertions, $q=0$ gives a constant empty product and $q=1$ gives a position-independent correlation by translation covariance. We now derive the directional estimate for $q\geq2$. Fix a compact set $K$ in the fixed-period cyclic patch. Since that patch is open, there is a compact neighborhood $K'$ of $K$ still contained in it. Compactness and continuity of the finitely many periodic coordinate-distance functions give a number $\rho>0$ such that every tuple in $K'$ has, for every pair of marks and every axis, distance at least $4\rho$ from the period lattice. Choose a centered-budget depth $m$ with $a_m<\rho$, and choose one common budget $A$ for the finite list consisting of the requested observables and their axis-reflected/conjugated versions. For sufficiently fine mesh $N$, the stencil fit conditions hold and the mesh is small enough for the following integer choices.

Move one marked site by $l\geq0$ lattice steps in an axis $\alpha$, where $l a_N\leq\rho$. For a negative displacement, reverse the endpoints and apply the same estimate along the reversed segment; the enlarged compact neighborhood keeps both paths admissible. Let $R=\rho/a_N$ be the corresponding radius in lattice units. The coordinate gaps in the cyclic order exceed the common brick width plus $2R$: the compact separation gives at least $4\rho/a_N$ lattice steps, whereas the brick has width at most $2a_m/a_N<2R$. Subtracting that width leaves free gaps greater than $2R$. The move changes two free gaps from $(H_0,H_{q-1})$ to $(H_0-l,H_{q-1}+l)$; all other gaps remain fixed, and total tour length is unchanged.

The centered root bound supplies the column estimate needed for this tour. Cut a marked observable at the transfer plane through its centered cell. Locality identifies it with a continuous one-dimensional slab observable, represented by a brick of width $w=2a_m/a_N$ in lattice units. The reflected cut gives a second slab function, obtained from the reflected token; both are admissible in the relevant home chamber and have centered root test at most $A$. The centered tiling has $4s$ cells along that axis. Put the $2s$ alternating original and reflected slab functions in the distinct even-indexed cells; their spacing is two cells, and the planted inequality applies to this injection. The planted inequality bounds the absolute expectation of their product by $A^{2s}$. Translation and reflection identify that product with the $s$-fold reflected column word of the brick. Thus the normalized column moment to power $\theta=1/(2s)$ is at most $A$.

We now justify the transfer estimate for the continuous slice kernels; they are not finite matrices. Fix a compact slice space $X$, finite measure $\mu$, continuous kernels $T,B_0,\ldots,B_{q-1}$, and the positive integer word length $L=2sw$. Choose finite Borel partitions $\{C_u^{(j)}\}_u$ of $X$ whose maximum cell diameter tends to zero, and choose a tag $\xi_u^{(j)}\in C_u^{(j)}$ in each nonempty cell. Form the sampled matrix
\[
 (K_j)_{uv}=\mu(C_u^{(j)})^{1/2}K(\xi_u^{(j)},\xi_v^{(j)})\mu(C_v^{(j)})^{1/2}.
\]
If $T$ is a positive-semidefinite kernel, $T_j$ is a positive-semidefinite matrix: it is the Gram matrix of the tagged kernel, conjugated by the diagonal matrix of square-root cell masses. For any fixed cyclic list $K_1,\ldots,K_r$, multiplication and taking the matrix trace gives the exact finite sum
\[
 \operatorname{Tr}((K_1)_j\cdots(K_r)_j)
 =\sum_{u_1,\ldots,u_r}\prod_{t=1}^{r}\mu(C_{u_t}^{(j)})K_t(\xi_{u_t}^{(j)},\xi_{u_{t+1}}^{(j)}),\qquad u_{r+1}=u_1.
\]
The integrand $\prod_tK_t(x_t,x_{t+1})$ is continuous on compact $X^r$. Since the cell diameters tend to zero, this tagged sum converges to its integral, the cyclic kernel integral $\operatorname{Cl}_\mu(K_1,\ldots,K_r)$. The same quadrature sequence works for the partition word, each shifted tour word, and each fixed replicated column word $((B_i)^*B_i)^s$, since these are finitely many continuous cyclic integrands. In particular, if $Z=\operatorname{Cl}_\mu(T,\ldots,T)>0$ uses $L$ copies, then $Z_j=\operatorname{Tr}(T_j^L)\to Z$ and $Z_j>0$ eventually. Define the adjoint kernel by $B_i^*(x,y)=\overline{B_i(y,x)}$. The column quantities $f_{i,j}$ converge to
\[
 f_i=\left|\operatorname{Cl}_\mu(B_i^*,B_i,\ldots,B_i^*,B_i)\right|^{1/(2s)},
\]
and the centered root bound just proved says $f_i\leq A Z^{1/(2s)}$.

For completeness, the finite-matrix estimate applied to these sampled matrices is an immediate Schatten--Hölder bound. For a finite matrix $A$, write $\|A\|_r=(\operatorname{Tr}|A|^r)^{1/r}$ for $1\leq r<\infty$ and use the operator norm for $r=\infty$. The two-factor Schatten Hölder inequality is $\|AB\|_r\leq\|A\|_u\|B\|_v$ when $1/r=1/u+1/v\leq1$. It follows from Tropp's generalized unitarily invariant norm Hölder inequality (\emph{Matrix Analysis}, 2022, Theorem 6.32, p.~52): take the unitarily invariant norm to be the Schatten $r$ norm and its conjugate powers to be $u/r$ and $v/r$. The endpoint cases follow by continuity of finite-dimensional Schatten norms as an exponent tends to infinity; when both tend to infinity this is the operator norm inequality. Iterating this two-factor statement with partial sums of reciprocal exponents, and finishing with $|\operatorname{Tr}A|\leq\|A\|_1$, gives $|\operatorname{Tr}(A_1\cdots A_k)|\leq\prod_j\|A_j\|_{r_j}$ whenever $\sum_jr_j^{-1}=1$. Put $Z_j=\operatorname{Tr}(T_j^L)$ and $\theta=1/(2s)$. For each insertion $B_{i,j}$ use exponent $2s$, so its Schatten norm is
\[
 f_{i,j}=\left(\operatorname{Tr}\big((B_{i,j}^*B_{i,j})^s\big)\right)^{1/(2s)}.
\]
For a heat factor $T_j^{h_i}$ use exponent $L/\operatorname{Re}h_i$; if $\operatorname{Re}h_i=0$, use $\infty$. The spectral theorem gives $\|T_j^{h_i}\|_{L/\operatorname{Re}h_i}=Z_j^{\operatorname{Re}h_i/L}$, and $\|T_j^{h_i}\|_\infty\leq1$ on the imaginary boundary after the zero-eigenvalue convention. The reciprocal exponents sum to one because $q/(2s)+\sum_i\operatorname{Re}h_i/L=1$. Hence
\[
 \left|\operatorname{Tr}\prod_i B_{i,j}T_j^{h_i}\right|
 \leq \left(\prod_i f_{i,j}\right)Z_j^{1-q\theta}.
\]
Now take $h_0=H_0-z$, $h_{q-1}=H_{q-1}+z$, and $h_i=H_i$ otherwise. Their sum is constant; each real part is nonnegative for $|z|\leq2R$ because every free gap exceeds $2R$. The matrix powers are entire in $z$ by spectral calculus, with the zero eigenspace assigned the constant-zero value. Thus this trace is an entire function $\Phi_j(z)$ satisfying
\[
 |\Phi_j(z)|\leq \left(\prod_i f_{i,j}\right)Z_j^{1-q\theta}.
\]
For any real $t\in[0,l]$, the radius-$R$ circle around $t$ is inside $|z|\leq2R$, so Cauchy's estimate gives $|\Phi_j'(t)|\leq (\prod_i f_{i,j})Z_j^{1-q\theta}/R$. Integrate from $0$ to $l$. Taking $j\to\infty$ in the resulting endpoint inequality is justified by the cyclic quadrature convergence above. It gives the continuous-kernel estimate
\[
 |F_{N,p,\mathbf O}(\mathbf x+l a_N e_{k,\alpha})-F_{N,p,\mathbf O}(\mathbf x)|
 \leq \frac{l}{R}\,\frac{(\prod_i f_i)Z^{1-q\theta}}{Z}
 \leq \frac{A^q}{\rho}\,l a_N.
\]
This is the required one-coordinate bound. It derives the continuous transfer estimate from finite sampled matrices and explicitly passes the traces, denominator, and all column norms to the limit; no finite-dimensional representation of the Wilson measure is assumed.

To compare arbitrary lattice tuples whose embeddings lie in $K$, change their $4q$ scalar coordinates one at a time. Intermediate tuples can lie outside $K$; by first choosing $K'$ as above and then taking $\delta$ smaller than its neighborhood margin, every intermediate tuple remains in $K'$. The same $A,\rho$ work at all stages. The sum of the $4q$ coordinate displacements is at most $4q$ times the sup distance, so
\[
 |F_{N,p,\mathbf O}(\mathbf x)-F_{N,p,\mathbf O}(\mathbf y)|
 \leq 4q\,\frac{A^q}{\rho}\,|\mathbf x-\mathbf y|.
\]
Choosing $\delta<\varepsilon\rho/(4qA^q)$, with any sufficiently small neighborhood margin if $A=0$, proves the fixed-period modulus.

For the expanding-period modulus, let $K$ be compact in the generic set, where each coordinate projection separates distinct points. Choose a compact neighborhood still inside that open set. Since its coordinate differences are bounded and the physical period tends to infinity, for every sufficiently large period all pairs are uniformly separated from every periodic image; the same positive $\rho$ works for all those periods. Choose the centered-budget depth $m$ with $a_m<\rho$. The expanding clause gives a common budget $A$ eventually in period, then eventually in mesh for each such period. The preceding coordinate telescope uses the same $\rho,A,\delta$ for all these periods, with the inner mesh threshold allowed to depend on the period. This is exactly the expanding modulus. Genericity, rather than absence of full point collisions alone, is needed here because all four coordinate gaps are used by the directional telescope.

We next prove both compact-bound assertions. The cases $q=0,1$ are immediate: the empty correlation is $1$, while for one observable the planted inequality with one occupied chamber bounds its expectation by its centered root test, at most $A$; translation covariance makes this uniform in its position. Hence assume $q\geq2$. Let $K$ be compact with no full collisions. Define for each pair $i\ne j$
\[
 d_{ij}(x)=\max_{\alpha=1,\ldots,4}|x_i^\alpha-x_j^\alpha|.
\]
Its minimum over the finite set of pairs and compact $K$ is a positive number. For sufficiently large periods, every coordinate difference has absolute value less than half the physical period; hence its distance to the period lattice is its absolute value. It follows that there is $\rho_0>0$ such that every pair in every tuple of $K$ is separated in at least one torus coordinate by $\rho_0$, uniformly on a period tail. The same conclusion for a fixed period holds under the weaker NoHit hypothesis: use instead the continuous positive function
\[
 d_{ij,p}(x)=\max_{\alpha=1,\ldots,4}\operatorname{dist}(x_i^\alpha-x_j^\alpha,P_p\mathbb Z).
\]
NoHit says precisely that $d_{ij,p}(x)>0$ for every pair; compactness and finiteness of the pair set give $\min_{x\in K,i\ne j}d_{ij,p}(x)>0$. This maximum over coordinates is essential: taking the minimum over coordinates would incorrectly rule out coordinatewise ties.

Here is the packing argument, with the two scales distinguished. Choose a coarse depth $m$ with $a_m<\rho_0$. Then choose a finer depth $k\geq m$ in the relevant centered-budget tail so that the budget at depth $k$ is available and
\[
 6q a_k<a_m.
\]
For expanding bounds use the expanding centered-budget clause; for a fixed-period bound use the fixed-$p$ clause. After fixing $k$, take the mesh index large enough that every fixed stencil fits inside a depth-$k$ chamber. The finite torus placement lemma supplies a translate of the coarse tiling $X$ of depth $m$ such that all $q$ marked sites lie in its margin of three fine-cell spans. The proof is a boundary count in each coordinate: one point forbids at most $6a_k$ of the coarse translation interval, so all $q$ forbid at most $6q a_k<a_m$; hence an allowed translate exists. Let $D_i$ be the coarse cell of $X$ containing $x_i$. These cells are distinct, since two sites in the same cell differ by less than $a_m<\rho_0$ in every coordinate, contradicting the some-coordinate separation.

The centered budgets are at the fine depth $k$, whereas the tuple will be placed in a coarse tiling $X$ of depth $m$. First choose the home site $o$ of $X$ on the fine center grid and in chamber zero of $X$. Match the entire coarse translated-copy orbit of the observable at $o$, one copy for each cell of $X$, injectively to cells of the fine centered tiling at $o$. The aligned scale divisibility supplies the matching; the required reflected and conjugated orientations are in the finite family for which the common budget was chosen. The fine planted inequality bounds the absolute expectation of the whole coarse copy product by $A^{|\mathrm{Cell}(X)|}$. By the definition of the root test, the root test of the observable at $o$ relative to the coarse tiling $X$ is at most $A$. This is the full replicated-orbit argument; bounding just the original tuple's $q$ entries would not prove a coarse root estimate.

For each marked point $x_i$, let $D_i$ be its assigned coarse cell. The tiling's pointed-site map chooses $y_i$ in the home chamber (the chamber based at cell zero), and reflection of the margin shows that $y_i$ has a margin of three fine-cell spans from the home-chamber boundary. The coarse cell orientation determines a reflected or conjugated token at $y_i$; signed-coordinate covariance identifies its copy in $D_i$ with the original observable at $x_i$. The finite list of all such oriented tokens was included in the common budget. We next move the coarse home point from $o$ to any such $y_i$. In each of four coordinates, bracket $y_i$ by two fine-center-grid points separated by at most one fine span $d$. The sixteen corner points lie in the margin of two spans. The fixed-stencil fit gives chamber interiority at every intermediate point on these four coordinate segments. For a fixed axis, let $u_j$ be the root test at successive integer translates of the full coarse copied product, and let $M_X$ be the number of cells. Translate the entire copied product by one lattice step. In each chamber the translated copy becomes one of the two neighboring translates, with the sign determined by that chamber's orientation. Reflection positivity applied to this translated product bounds its moment by the planted product of the root tests of those neighboring translates. Pairing cells by the involution that reverses the orientation along this axis gives exactly $M_X/2$ copies of each neighbor. Thus
\[
 u_j^{M_X}\leq u_{j-1}^{M_X/2}u_{j+1}^{M_X/2},
\]
and raising both sides to the power $2/M_X$ gives $u_j^2\leq u_{j-1}u_{j+1}$. If an interior value exceeds both endpoint bounds, choose a maximum $u_j$ on the finite interval; it is positive. Both neighbors are at most $u_j$, so the inequality forces both to equal $u_j$. Repeating this at the adjacent interior maxima propagates that value to an endpoint, contradicting the endpoint bound. This proves the finite-interval maximum principle, including possible zeros elsewhere. Repeatedly applying this maximum principle along the four segments bounds the root at $y_i$ by the maximum of the sixteen corner roots, all at most $A$. This proves the coarse root estimate at each needed marked home point, with chamber interiority there; covariance transfers that bound to the coarse insertion associated with $x_i$.

Now apply the planted inequality in the coarse tiling $X$ to the injection $i\mapsto D_i$. It bounds the tuple expectation by the product of the $q$ coarse root tests, hence by $A^q$. All choices were uniform over tuples in $K$, so this is the expanding compact bound; for fixed $p$ it gives the fixed-period compact NoHit bound. The former uses compactness away from full collisions and the period tail; the latter uses only NoHit at that fixed period and permits coordinate ties. Neither gives control at a torus-coincidence configuration, nor smoothness through a collision diagonal.
\end{proof}

\begin{remark}[What remains upstream]\label{ex-centered-open}
The theorem assumes the centered budget. The exact upstream estimate needed is: for each grade $d$, norm bound $R>0$, and exponent slack $\eta>0$, there is $A\geq0$ such that, for every depth $m$ beyond the reference depth for grade $d$, eventually in volume index $N$, uniformly for every period index $p$ and every label $v$ of grade $d$ with $\|v\|\leq R$, the centered root test is at most $A a_m^{-d-\eta}$, with representation and chamber-interiority. The same volume tail works for all $p$; restricting either to a fixed period or to the cofinal period schedule gives the two budget regimes used above. No proof of this Wilson estimate is contained here. In particular, the finite-volume cell-packing and transfer arguments above cannot supply its decay in $a_m$, nor can they supply the required uniformity in labels of bounded norm. This is the next analytic input to establish from the renormalisation-group trajectory.
\end{remark}
